%% ---- 5. Configuração Experimental -------------------------------------------
\section{Configuração Experimental}

\paragraph{Dados e modelos.}
O conjunto principal é o MovieLens-20M~\cite{harper2015movielens} (138\,493
usuários, 18\,345 filmes), que tem estrutura sequencial forte e históricos
longos~\cite{klenitskiy2026patterns}, o que mantém $2^n$ fora do alcance da
enumeração. O segundo é o Zvuk~\cite{shevchenko2024zvuk}, um registro de
audições de música, do qual usamos 100\,000 usuários e as 50\,000 faixas mais
ouvidas. Em ambos, o recomendador é um SASRec~\cite{kang2018sasrec} treinado
com entropia cruzada sobre todos os itens~\cite{klenitskiy2023turning}, com
comprimento máximo 128 e divisão \emph{leave-one-out}. O pré-processamento, a
busca de hiperparâmetros no Zvuk e as verificações de que os modelos usam a
ordem do histórico estão no Apêndice~\ref{ap:treino}. Os dois modelos atingem
NDCG@10 de \val{train.ml.test.ndcg} (MovieLens) e \val{train.zv.test.ndcg}
(Zvuk) no teste, com ranking completo~\cite{krichene2020sampled}.

\paragraph{Instâncias.}
Uma instância é um par (usuário, item-alvo). O histórico $s$ são os itens do
usuário anteriores ao item de teste, limitados aos 128 mais recentes. Os
usuários têm $n\ge30$ e vêm de três estratos de comprimento ($30$--$49$,
$50$--$99$, $100$--$128$). O alvo $v$ é sorteado uniformemente entre os 10
primeiros candidatos. Mantemos só instâncias em que apagar o histórico inteiro
dá $\hat N<0$, isto é, em que existe ao menos uma explicação válida; nenhuma
instância sorteada foi descartada por esse critério. No MovieLens usamos três
conjuntos de usuários disjuntos: o \textbf{principal} (33 por estrato, 99
instâncias), o de \textbf{ajuste} (10 por estrato), usado só para escolher
parâmetros, e um conjunto \textbf{novo} (99 instâncias), sorteado depois dos
resultados principais para verificar que eles não dependem da amostra. No Zvuk
usamos 45 instâncias (15 por estrato).

\paragraph{Orçamento.}
Cada execução de cada método recebe $10^5$ passadas do SASRec. Com $M=10$,
isso são 5\,000 máscaras (10\,000 para o NSGA-II só $(N,C)$, que não estima
$\hat S$); máscaras repetidas não são cobradas. A reavaliação
com $M'=100$ fundos fica fora do orçamento. Para o GECE, o custo é o número de
sequências que ele consulta. Com $M=10$ a ordenação das máscaras concorda com a
de 200 fundos (Spearman de \val{bench.ml.n} em $\hat N$ e \val{bench.ml.s} em
$\hat S$ no MovieLens, \val{bench.zv.n} e \val{bench.zv.s} no Zvuk).

\paragraph{Métricas.}
A métrica principal é o hipervolume~\cite{zitzler1999spea} das máscaras com no
máximo 10 posições (\HVp), a região de explicações que uma pessoa leria; a
secundária é o hipervolume de todas as máscaras (HV). A proposta do projeto
previa só o HV; o \HVp{} entrou na análise antes dos relatórios da comparação
principal, e reportamos sempre os dois. Em cada instância, os três
objetivos da frente reavaliada são normalizados para $[0,1]$ pelos menores e
maiores valores entre todas as frentes comparadas naquela instância (no \HVp,
só entre as máscaras pequenas), com ponto de referência $1{,}1$ em cada
objetivo; o hipervolume vai, então, de 0 a $1{,}1^3\approx1{,}33$. Como a
normalização depende do conjunto de frentes comparadas, valores de relatórios
diferentes não são comparáveis entre si. O IGD+~\cite{ishibuchi2015igdplus}, a
maior $\widehat{\PN}$ e a maior $\widehat{\PS}$ de cada frente ficam no
Apêndice~\ref{ap:metricas}: o IGD+ ordena os métodos como o \HVp{}, e as duas
probabilidades quase não distinguem os métodos de busca. O SMS-EMOA seleciona
pela contribuição ao hipervolume, a mesma grandeza usada para avaliá-lo; seu
HV de todas as máscaras, porém, fica abaixo do do NSGA-II, de modo que a
vantagem não é automática.

\paragraph{Protocolo estatístico.}
Métodos estocásticos rodam com 5 sementes no conjunto principal e 3 nos
demais. A unidade de análise é a instância: fazemos a média das sementes dentro
de cada instância e comparamos métodos pelo teste de Wilcoxon pareado
bilateral, com correção de Holm em duas famílias: o melhor método contra cada
um dos demais, e um conjunto de pares definidos depois das execuções para
responder a perguntas específicas (Tabela~\ref{tab:pares}). Para cada par
reportamos a diferença média com intervalo de 95\% por \emph{bootstrap} sobre
instâncias e a correlação bisserial de postos pareada ($-1$ a $+1$).

\paragraph{Comparação planejada e experimentos posteriores.}
\label{sec:posthoc}
A \textbf{comparação planejada}, definida na proposta do projeto, opõe o
NSGA-II ao AG escalarizado, ao guloso, ao ranking por ablação, ao GECE e às
heurísticas (Seção~\ref{sec:planejada}). Os demais experimentos foram
decididos depois dos primeiros resultados, em resposta a eles, e estão
marcados com $\dagger$: os outros três algoritmos
multiobjetivo, o NSGA-II semeado, o NSGA-II só $(N,C)$ e a reavaliação de mais
candidatas suas, o Shapley e as configurações ajustadas do AG escalarizado
(Seção~\ref{sec:posteriores}).
Todos usam os mesmos fundos, instâncias e orçamento, e cada parâmetro novo foi
escolhido nos usuários de ajuste. Também o conjunto novo do MovieLens e o Zvuk
foram acrescentados depois dos primeiros resultados, para verificar que as
conclusões não dependem da amostra nem do modelo (Seção~\ref{sec:robustez}).
Nenhuma decisão de método ou parâmetro usou esses dados. O NSGA-II semeado e as
conclusões da comparação planejada foram obtidos antes deles; os outros três
algoritmos multiobjetivo rodaram na mesma leva que eles, de modo que, para
esses, os dados novos são uma replicação simultânea, e não uma confirmação
posterior.

\paragraph{Execução.}
Um escalonador avança muitas execuções em paralelo e envia as máscaras de
todas elas em um único lote à GPU (até \val{bench.ml.speed} passadas por
segundo em uma RTX~A4000 no MovieLens). Todos os experimentos, incluindo o
treino dos recomendadores, custaram menos de seis dólares em GPUs alugadas.
